% =====================================================================
%  BÁO CÁO GIỮA KỲ — KHÓA LUẬN TỐT NGHIỆP
%  Thiết kế CNN Accelerator cho phân loại rối loạn nhịp tim ECG
%  trên FPGA Intel Cyclone V
%  SV: Lê Minh Đức (22200034)
%
%  TRÊN OVERLEAF: chọn compiler = XeLaTeX (Menu > Settings > Compiler).
%    Dùng font chuẩn LaTeX = Latin Modern (bản Unicode của Computer Modern),
%    nạp qua gói fontspec/unicode-math — có sẵn trên Overleaf, hỗ trợ tiếng Việt.
% =====================================================================
\documentclass[12pt,a4paper]{article}

\usepackage{fontspec}
\usepackage{polyglossia}
\setdefaultlanguage{vietnamese}
% Latin Modern = font chuẩn của LaTeX (Computer Modern hiện đại, có Unicode)
\setmainfont{Latin Modern Roman}
\setmonofont{Latin Modern Mono}  % monospace cho code/diagram block

\usepackage{amsmath,amssymb}
\usepackage{geometry}
\geometry{a4paper,top=2.5cm,bottom=2.5cm,left=3cm,right=2cm}
\usepackage{booktabs}
\usepackage{array}
\usepackage{tabularx}
\usepackage{longtable}
\usepackage{xcolor}
\usepackage{listings}
\usepackage{titlesec}
\usepackage{enumitem}
\usepackage[hidelinks]{hyperref}
\usepackage{caption}

% ----- màu & tiêu đề -----
\definecolor{accent}{RGB}{31,58,95}
\titleformat{\section}{\Large\bfseries\color{accent}}{\thesection.}{0.5em}{}
\titleformat{\subsection}{\large\bfseries\color{accent}}{\thesubsection.}{0.5em}{}
\setlength{\parskip}{4pt}
\setlength{\parindent}{0pt}
\renewcommand{\arraystretch}{1.25}
\captionsetup{font=it,labelsep=period}
\setlist[itemize]{leftmargin=1.4em,itemsep=2pt,topsep=2pt}

% ----- style cho code/diagram block -----
\lstset{
  basicstyle=\ttfamily\small,
  breaklines=true,
  columns=fullflexible,
  frame=none,
  xleftmargin=1.2em,
  keepspaces=true,
}

\begin{document}

% =====================================================================
%  TRANG BÌA
% =====================================================================
\begin{center}
  \vspace*{1.5cm}
  {\LARGE\bfseries\color{accent} BÁO CÁO GIỮA KỲ}\\[4pt]
  {\LARGE\bfseries\color{accent} KHÓA LUẬN TỐT NGHIỆP}\\[18pt]
  {\large\bfseries Thiết kế CNN Accelerator cho phân loại rối loạn nhịp tim ECG\\
   trên FPGA Intel Cyclone V}\\[2cm]
\end{center}

\begin{flushleft}
  \hspace{2cm}\textbf{Sinh viên thực hiện:}\quad Lê Minh Đức (22200034)\\[6pt]
  \hspace{2cm}\textbf{Giảng viên hướng dẫn:}\quad ThS.\ Nguyễn Thị Thiên Trang\\[2pt]
  \hspace{2cm}\hphantom{\textbf{Giảng viên hướng dẫn:}}\quad ThS.\ Mã Khải Minh\\[6pt]
  \hspace{2cm}\textbf{Ngày báo cáo:}\quad 01/07/2026\\
\end{flushleft}

\vfill

\clearpage

% =====================================================================
%  1. TÓM TẮT
% =====================================================================
\section{Tóm tắt}

Khóa luận thiết kế và triển khai một lõi IP CNN Accelerator nhẹ (lightweight) trên FPGA
Intel Cyclone V nhằm phân loại rối loạn nhịp tim từ tín hiệu điện tim (ECG) đơn đạo trình
(lead II). Dự án gồm hai phần gắn kết chặt chẽ: (i) một pipeline phần mềm PyTorch thực hiện
huấn luyện --- cắt tỉa (pruning) --- lượng tử hóa INT8 lũy thừa hai (power-of-two), và (ii)
một lõi phần cứng Verilog đã được kiểm chứng bit-exact với mô hình phần mềm và chạy thực tế
trên kit DE10-Standard.

\textbf{Các kết quả chính tính đến thời điểm báo cáo:}
\begin{itemize}
  \item \textbf{Độ chính xác:} $\sim$94\% / macro-F1 $\sim$0,93 trên tập Chapman (4 lớp
        AFIB/GSVT/SB/SR, chia patient-independent 70/15/15), với mô hình chỉ 654 tham số.
  \item \textbf{Bit-exact:} 21 điểm kiểm tra (checkpoint) mỗi mẫu khớp tuyệt đối giữa Python
        và RTL --- $\max|\text{diff}| = 0$ LSB trên 15.312 phần tử so sánh. Lượng tử hóa INT8
        không gây mất mát so với float32 (cùng $\sim$94\%).
  \item \textbf{Phần cứng:} một inference hoàn tất trong 5.216 chu kỳ $\approx$ 52,16 µs @
        100 MHz (xác định, không biến thiên), dùng 2.201 ALM (5\%) và 28 DSP (25\%) của thiết
        bị 5CSXFC6D6F31C6; Fmax = 104,85 MHz.
  \item \textbf{Demo trên board:} DE10-Standard qua cầu JTAG-to-Avalon đạt 94,27\%
        (1004/1065), tái lập kết quả mô phỏng.
  \item \textbf{Cross-dataset:} mô hình huấn luyện trên bộ Chapman mở rộng (Chapman/Ningbo)
        đánh giá zero-shot trên Georgia (PhysioNet 2020, 12-lead, 5459 record) đạt 93,00\% /
        macro-F1 0,9151 / macro-AUC 0,9580 (INT8 bit-exact), chứng minh lõi tổng quát hóa thật
        ra ngoài dữ liệu huấn luyện.
\end{itemize}

% =====================================================================
%  2. MỤC TIÊU
% =====================================================================
\section{Mục tiêu đề tài và bài toán}

Rối loạn nhịp tim (cardiac arrhythmia) là một trong những nguyên nhân hàng đầu gây đột tử
tim. Việc phát hiện sớm bằng thiết bị đeo (wearable) giám sát liên tục có giá trị lâm sàng
cao. Mạng nơ-ron tích chập một chiều (1D-CNN) trên tín hiệu ECG thô đạt độ chính xác tốt,
tuy nhiên vi điều khiển (MCU) thiếu thông lượng cho inference liên tục, còn GPU biên (edge
GPU) lại không phù hợp về công suất và kích thước cho thiết bị đeo. FPGA Cyclone V là điểm
cân bằng hợp lý: đủ song song hóa cho CNN nhỏ, công suất thấp, và có thể nhúng soft-logic
hoặc cầu host để điều khiển datapath.

\textbf{Mục tiêu thiết kế cụ thể của đề tài:}
\begin{itemize}
  \item Thiết kế một mô hình CNN 1D nhỏ gọn, đạt độ chính xác $\geq \sim$94\% trên 4 lớp nhịp
        Chapman, trên cơ sở tham khảo và tối ưu kiến trúc của bài báo Liu (2023).
  \item Lượng tử hóa INT8 với khâu rescale không cần multiplier (power-of-two) nhằm giảm tài
        nguyên DSP và năng lượng tiêu thụ.
  \item Triển khai lõi RTL được kiểm chứng bit-exact với mô hình phần mềm (không phải chỉ là
        ``xấp xỉ INT8'' như đa số công trình trước).
  \item Kiểm chứng chức năng thiết kế bằng mô phỏng và triển khai, đo đạc thực tế trên board
        DE10-Standard.
  \item Kiểm tra khả năng tổng quát hóa (cross-dataset) trên dữ liệu Georgia 12-lead để đánh
        giá độ bền của mô hình ngoài phân phối huấn luyện.
\end{itemize}

% =====================================================================
%  3. SOFTWARE
% =====================================================================
\section{Đề xuất mô hình và tối ưu pipeline (Software)}

\subsection{Kiến trúc mô hình đề xuất}

Mô hình ECG\_1DCNN (bản đã pruned) gồm 4 lớp tích chập một chiều (kênh 4-4-8-8, kernel K=5,
padding=2), mỗi lớp theo sau bởi MaxPool stride 5; tiếp đến là Global Average Pooling (GAP)
$\rightarrow$ lớp Fully-Connected (8$\rightarrow$4) $\rightarrow$ argmax. Kiến trúc kế thừa
tinh thần ``fully-mapped 1D-CNN'' nhưng được thu nhỏ mạnh (số kênh là lũy thừa 2 để thuận
lợi cho phần cứng) và cắt tỉa kênh để giảm tham số xuống còn 654 --- cực nhỏ, hướng tới
thiết bị đeo.

\begin{table}[h]
\centering
\caption{Cấu trúc mô hình ECG\_1DCNN pruned (654 tham số).}
\begin{tabular}{lcccccrr}
\toprule
\textbf{Lớp} & \textbf{In\_ch} & \textbf{Out\_ch} & \textbf{K} & \textbf{Pool} & \textbf{ReLU} & \textbf{In\_len} & \textbf{Out\_len}\\
\midrule
Conv1 & 1 & 4 & 5 & /5 & Không & 2500 & 500\\
Conv2 & 4 & 4 & 5 & /5 & Không & 500 & 100\\
Conv3 & 4 & 8 & 5 & /5 & Không & 100 & 20\\
Conv4 & 8 & 8 & 5 & /5 & Có & 20 & 4\\
GAP & 8 & 8 & --- & /4 & --- & 4 & 1\\
FC & 8 & 4 & --- & --- & --- & 1 & 1\\
\bottomrule
\end{tabular}
\end{table}

\begin{itemize}
  \item ReLU chỉ đặt sau Conv4 --- nhằm giữ lại đặc trưng âm của tín hiệu ECG ở Conv1--Conv3.
  \item Đầu vào: 2500 mẫu INT8 (5 giây @ 500 Hz, lead II). Đầu ra: 4 lớp AFIB / GSVT / SB / SR.
\end{itemize}

\subsection{Thuật toán suy luận của mô hình CNN}

Lan truyền tiến (forward) của mô hình là chuỗi bốn khối tích chập--gộp giống nhau, kết thúc
bằng gộp toàn cục và một lớp tuyến tính. Mỗi lớp tích chập 1D thực chất là phép tương quan
chéo (cross-correlation) trượt một kernel độ dài K=5 dọc trục thời gian; với mỗi vị trí đầu
ra $t$ và mỗi kênh ra $j$:
\begin{equation}
  y_j[t] = b_j + \sum_{c=1}^{C_{in}} \sum_{k=0}^{K-1} w_{j,c}[k]\;\cdot\;x_c\,[\,t + k - P\,]
\end{equation}
trong đó $P=2$ là padding (đệm 0 hai biên để giữ độ dài), $C_{in}$ là số kênh vào. Vì padding
``same'', độ dài chuỗi không đổi qua tích chập; việc giảm chiều hoàn toàn do MaxPool đảm
nhiệm. MaxPool 1D với cửa sổ và stride đều bằng 5 lấy giá trị lớn nhất trên từng cụm 5 mẫu
liên tiếp, giảm độ dài đúng 5 lần mỗi lớp (2500$\rightarrow$500$\rightarrow$100$\rightarrow$20$\rightarrow$4):
\begin{equation}
  p_j[t] = \max_{0 \le k < 5}\; y_j[\,5t + k\,]
\end{equation}
Sau bốn khối, mỗi kênh còn 4 mẫu. Global Average Pooling (GAP) nén mỗi kênh thành một số bằng
trung bình trên trục thời gian, tạo vector đặc trưng độ dài bằng số kênh (8):
\begin{equation}
  g_j = \frac{1}{4}\sum_{t=0}^{3} p_j[t]
\end{equation}
Vector đặc trưng đi qua lớp Fully-Connected (8$\rightarrow$4) cho 4 logit, rồi argmax chọn
lớp có logit lớn nhất làm nhãn dự đoán:
\begin{equation}
  z_i = \sum_{j} W_{i,j}\,g_j + b_i, \qquad \hat{y} = \arg\max_i\, z_i
\end{equation}
\begin{itemize}
  \item ReLU chỉ đặt trước FC (sau Conv4) là một lựa chọn có chủ đích: nhịp tim mang thông tin
        ở cả phần âm lẫn dương của tín hiệu, nên Conv1--3 không ReLU để không cắt mất biên độ
        âm; ReLU ở Conv4 tạo phi tuyến tính cuối cùng cho bộ phân loại tuyến tính phía sau.
\end{itemize}

\subsection{Hàm mất mát, lan truyền ngược và bộ tối ưu}

Mô hình được huấn luyện theo học có giám sát. Đầu ra 4 logit được chuẩn hóa thành phân phối
xác suất bằng softmax, và hàm mất mát là cross-entropy giữa phân phối dự đoán và nhãn thật
(one-hot). Với một mẫu có nhãn đúng là lớp $c$:
\begin{equation}
  \mathcal{L} = -\log\!\left( \frac{e^{z_c}}{\sum_{i} e^{z_i}} \right) = -z_c + \log\!\sum_{i} e^{z_i}
\end{equation}
Trong PyTorch, hai bước softmax và cross-entropy được gộp trong \texttt{nn.CrossEntropyLoss}
(nhận thẳng logit, ổn định số học). Mất mát của cả batch là trung bình theo mẫu.

Quá trình lan truyền ngược (backpropagation) tính gradient của hàm mất mát theo mọi tham số
bằng quy tắc chuỗi, lan ngược từ lớp FC qua GAP, các khối Conv--Pool về tới Conv1. Một bước
huấn luyện gồm bốn thao tác lặp lại trên từng batch:

\begin{lstlisting}
for each batch (ecg, label):
    optimizer.zero_grad()        # xoa gradient tich luy
    logits = model(ecg)          # lan truyen tien (forward)
    loss   = CrossEntropy(logits, label)
    loss.backward()              # lan truyen nguoc: tinh dL/dtheta
    optimizer.step()             # cap nhat tham so theta
\end{lstlisting}

Bộ tối ưu là Adam (Adaptive Moment Estimation) --- kết hợp động lượng (moment bậc nhất) và
tỉ lệ học thích nghi theo từng tham số (moment bậc hai), hội tụ nhanh và ít nhạy với việc
chọn tỉ lệ học hơn SGD thuần. Với mỗi tham số $\theta$, Adam duy trì hai ước lượng trung bình
trượt của gradient $g$ và bình phương gradient:
\begin{equation}
  m_t = \beta_1 m_{t-1} + (1-\beta_1)\,g_t, \qquad v_t = \beta_2 v_{t-1} + (1-\beta_2)\,g_t^{2}
\end{equation}
\begin{equation}
  \hat{m}_t = \frac{m_t}{1-\beta_1^{t}}, \quad \hat{v}_t = \frac{v_t}{1-\beta_2^{t}}, \qquad
  \theta_t = \theta_{t-1} - \eta\,\frac{\hat{m}_t}{\sqrt{\hat{v}_t}+\epsilon}
\end{equation}

\textbf{Cấu hình huấn luyện thực tế:}
\begin{itemize}
  \item Hàm mất mát: CrossEntropyLoss; bộ tối ưu: Adam, tỉ lệ học khởi đầu $\eta = 10^{-3}$.
  \item Lịch tỉ lệ học: MultiStepLR giảm $\times 0,1$ tại giữa quá trình (epoch 50)
        $\rightarrow$ tinh chỉnh ở $10^{-4}$.
  \item 100 epoch, batch size 128; chọn checkpoint tốt nhất theo val loss thấp nhất
        (early-stopping ngầm theo validation).
  \item Cùng công thức tối ưu (Adam + CrossEntropy) được tái sử dụng ở cả hai giai đoạn sau:
        fine-tune sau pruning và huấn luyện QAT --- chỉ khác tỉ lệ học và số epoch.
\end{itemize}

\subsection{Các thuật toán tối ưu mô hình (nén \& lượng tử hóa)}

\textbf{Ba kỹ thuật tối ưu được áp dụng tuần tự để đưa mô hình từ một CNN float32 đầy đủ về
một lõi INT8 cực nhỏ, phù hợp phần cứng mà giữ độ chính xác.}

\textbf{(a) Cắt tỉa kênh có cấu trúc (structured channel pruning).}
Mô hình gốc (kênh 4-8-8-16) được cắt tỉa xuống (4-4-8-8) bằng cách loại bỏ trọn các bộ lọc
(output channel) ít quan trọng nhất --- cắt tỉa ``có cấu trúc'' nên kết quả vẫn là một CNN
dày đặc, chạy thẳng trên phần cứng mà không cần hỗ trợ ma trận thưa. Độ quan trọng của mỗi
kênh được xếp hạng theo hai tiêu chí:
\begin{itemize}
  \item Chuẩn L1 của bộ lọc: kênh có tổng trị tuyệt đối trọng số nhỏ $\rightarrow$ đóng góp
        ít $\rightarrow$ cắt trước.
\end{itemize}
\begin{equation}
  I_j^{L1} = \sum_{c,k} \left| w_{j,c}[k] \right|
\end{equation}
\begin{itemize}
  \item Điểm Taylor bậc nhất (Molchanov et al., 2017): ước lượng mức tăng hàm mất mát nếu xóa
        kênh $j$ bằng tích của activation và gradient, lấy trung bình trên dữ liệu --- bám sát
        ảnh hưởng thật tới độ chính xác hơn chuẩn L1 thuần.
\end{itemize}
\begin{equation}
  I_j^{Taylor} = \left| \frac{1}{N}\sum_{n}\; a_j^{(n)} \cdot \frac{\partial \mathcal{L}}{\partial a_j^{(n)}} \right|
\end{equation}
Sau khi cắt, mô hình được tinh chỉnh lại (fine-tune) hai pha để phục hồi độ chính xác: pha 1
(30 epoch, lr=$10^{-3}$) khôi phục nhanh, pha 2 (20 epoch, lr=$10^{-4}$) tinh chỉnh. Kết quả:
tham số giảm từ $\sim$1244 xuống 654 mà accuracy gần như không đổi ($\sim$94\%).

\textbf{(b) Lượng tử hóa nhận biết huấn luyện (QAT) với STE và scale EMA.}
Thay vì chỉ lượng tử hóa sau huấn luyện, mô hình được huấn luyện trực tiếp dưới điều kiện
lượng tử hóa ``giả'' (fake quantization): trong forward, mỗi tensor trọng số/activation bị
lượng tử hóa rồi giải-lượng-tử ngay (quantize--dequantize) để mạng ``nhìn thấy'' sai số lượng
tử ngay khi học:
\begin{equation}
  x_q = \mathrm{round}\!\left(\frac{x}{s}\right)\cdot s, \quad \text{clamp về } [-127, 127]
\end{equation}
\begin{itemize}
  \item Straight-Through Estimator (STE): phép round không khả vi, nên ở backward gradient
        được truyền thẳng qua (đạo hàm xấp xỉ bằng 1), cho phép huấn luyện qua khâu lượng tử
        hóa.
\end{itemize}
\begin{equation}
  \text{forward: } x_q;\qquad \text{backward: } \frac{\partial x_q}{\partial x} \approx 1
\end{equation}
\begin{itemize}
  \item Scale học bằng trung bình trượt mũ (EMA) của biên độ lớn nhất qua các batch, ổn định
        hơn lấy max tức thời; sau giai đoạn warmup, scale được ``đóng băng'' để hội tụ ổn định.
\end{itemize}

\textbf{(c) Lượng tử hóa power-of-two} (chi tiết công thức ở mục 3.5).
Khác biệt then chốt so với QAT thông thường: scale của mỗi lớp bị ràng buộc là một lũy thừa
của hai, biến phép rescale đắt đỏ (nhân với số thực + dịch) thành chỉ một phép dịch bit + cộng
bù làm tròn. Đây chính là cầu nối giữa tối ưu phần mềm và tối ưu phần cứng: lựa chọn ở khâu
huấn luyện trực tiếp xóa bỏ multiplier ở khâu rescale trên FPGA.

\subsection{Lượng tử hóa power-of-two QAT --- công thức}

Trọng số, activation và bias được lượng tử hóa INT8 với scale là lũy thừa hai theo từng lớp.
Số bit dịch của mỗi lớp được chọn theo:
\begin{equation}
  nb = \left\lfloor \log_2\!\left(\frac{127}{|x|_{\max}}\right) \right\rfloor
\end{equation}
Giá trị cụ thể cho từng lớp (Conv1..Conv4, FC):
\begin{equation}
  nb = \{8,6,6,7,0\}, \quad w\_shift = \{6,6,6,7,8\}, \quad input\_shift = 2
\end{equation}
Phép rescale ở đầu ra mỗi lớp dùng làm tròn nửa lên (round-half-up) rồi bão hòa về dải INT8:
\begin{equation}
  out = \mathrm{clamp}\!\left( \left\lfloor \frac{acc + 2^{\,nb-1}}{2^{\,nb}} \right\rfloor,\ -127,\ 127 \right)
\end{equation}
Tương đương với phép dịch bit số học có cộng bù làm tròn:
\begin{equation}
  round\_half\_up(x) = (x + 2^{\,nb-1}) \gg nb
\end{equation}
Bias được scale theo cùng hệ số và lưu INT32 little-endian:
\begin{equation}
  bias\_scaled = \mathrm{round}\!\left( b_{float} \cdot 2^{\,nb} \right)
\end{equation}
Vì mọi scale là lũy thừa hai, khâu rescale chỉ cần dịch bit + cộng, tức 0 DSP --- trong khi
lượng tử general-scale cần 1 multiplier cho mỗi điểm rescale. Đây là đòn bẩy tiết kiệm tài
nguyên và năng lượng then chốt của thiết kế.

\subsection{Kết quả lượng tử}

Để định lượng đánh đổi của power-of-two so với general-scale, dự án thực hiện một ma trận
ablation (single-run, seed=42; cột DSP là số multiplier cần cho riêng khâu rescale):

\begin{table}[h]
\centering
\caption{Ablation lượng tử hóa trên Chapman.}
\begin{tabular}{llcrrr}
\toprule
\textbf{Variant} & \textbf{Scale} & \textbf{Train} & \textbf{Acc \%} & \textbf{F1} & \textbf{DSP rescale}\\
\midrule
Float32 baseline & --- & --- & 94,65 & 0,9402 & ---\\
PTQ power-of-2 & $2^{nb}$ & none & 94,08 & 0,9338 & 0\\
PTQ general & absmax/127 & none & 94,46 & 0,9380 & 4\\
QAT power-of-2 (ours) & $2^{nb}$ & fake-quant & 94,37 & 0,9364 & 0\\
QAT general & absmax/127 & fake-quant & 94,65 & 0,9398 & 4\\
\bottomrule
\end{tabular}
\end{table}

% =====================================================================
%  4. HARDWARE
% =====================================================================
\section{Thiết kế phần cứng dựa trên thuật toán (Hardware)}

\subsection{Kiến trúc tổng thể}

\begin{lstlisting}
Avalon-MM  ->  Input SRAM (2500 x 8b)
                   |  (MUX cho Conv1)
         Conv-Pool Engine (8 CP block song song,
         K=5 pad=2, MaxPool /5)
                   |  Ping-Pong SRAM (feature map lien lop)
         GAP / FC / Argmax  ->  result[1:0]  (0..3)
\end{lstlisting}

Một FSM điều khiển phẳng 8 trạng thái tuần tự hóa 4 lớp conv cùng đuôi GAP/FC/Argmax. Module
\texttt{avalon\_slave.v} đóng vai bus adapter, dùng chung cho cả luồng synthesis (virtual-pin)
lẫn luồng on-board (JTAG-to-Avalon).

\subsection{CP block --- pipeline 5 tầng, tách 3 submodule}
\begin{itemize}
  \item MAC (S1--S4): 5 multiplier INT8$\times$INT8 (DSP18) $\rightarrow$ adder tree 3 tầng
        $\rightarrow$ tree\_out.
  \item Accumulate + rescale (S5--S8): cộng dồn partial sum theo in\_ch; gập (fold) bias +
        round-add vào giá trị khởi tạo accumulator để cắt khỏi critical path (kết quả số học
        không đổi); clamp $[-127,127]$; ReLU cho Conv4.
  \item Pool (S9): comparator max trượt cho MaxPool K=5, stride 5.
\end{itemize}
Độ sâu pipeline từ MUX cửa sổ vào tới thời điểm cập nhật accumulator là đúng 5 chu kỳ; delay
chain trong controller phải khớp con số này --- một lỗi off-by-one ở đây từng là bug kiểm
chứng chính và đã được sửa.

\subsection{Dataflow streaming gập (folded)}

Thay vì map toàn bộ mỗi lớp ra phần cứng riêng (fully-mapped như competitor Liu 2023),
Conv-Pool Engine time-multiplex 8 CP block trên các vị trí output và stream đầu vào qua
shift-register. Đây là đòn bẩy diện tích chính: tái sử dụng một bộ PE cho mọi vị trí của một
lớp, đổi độ trễ tăng (nhưng xác định) lấy mức giảm logic/register đáng kể.

\subsection{Số liệu phần cứng}

Thiết bị Cyclone V 5CSXFC6D6F31C6, Fmax đọc ở góc Slow 1100mV/85°C:

\begin{table}[h]
\centering
\caption{Tài nguyên \& timing trên Cyclone V.}
\begin{tabular}{lrr}
\toprule
\textbf{Metric} & \textbf{Baseline (FF-ROM)} & \textbf{+ weight-RAM reload}\\
\midrule
ALM & 2.201 (5\%) & 2.820 (7\%)\\
DSP & 28 (25\%) & 28 (25\%)\\
Registers & 3.177 & 4.852\\
M10K & 20 (4\%) & 28 (5\%)\\
Fmax (standalone) & 104,85 MHz & 108,94 MHz\\
Latency & 5.216 cy (52,16 µs) & 5.216 cy (52,16 µs)\\
Throughput & $\sim$19.200 inf/s & $\sim$19.200 inf/s\\
\bottomrule
\end{tabular}
\end{table}

\begin{itemize}
  \item Overhead của weight-reload (+619 ALM, +8 M10K) đến từ logic nạp-được (read-address
        adder, 8-way write decode, lắp ráp 40-bit), không phải từ bộ nhớ (M10K = 0 ALM).
\end{itemize}

% =====================================================================
%  5. VERIFICATION
% =====================================================================
\section{Kiểm định chức năng thiết kế (Verification)}

Hợp đồng kiểm chứng của dự án là bit-exactness giữa Python và RTL. Mô hình golden Python tái
lập đúng trình tự số học của RTL:
\begin{equation}
\begin{aligned}
  acc_{32} \;\rightarrow\; &+\,bias\_scaled \;\rightarrow\; +\,2^{\,nb-1} \;\rightarrow\; \gg nb\\
  \;\rightarrow\; &\mathrm{clamp}[-127,127] \;\rightarrow\; \mathrm{ReLU}(\text{Conv4}) \;\rightarrow\; \mathrm{MaxPool}
\end{aligned}
\end{equation}
Trong đó GAP dùng phép chia nguyên (floor) và FC giữ nguyên logit INT32 ($nb=0$) đưa vào
argmax:
\begin{equation}
  GAP = \left\lfloor \frac{sum}{4} \right\rfloor = sum \gg 2
\end{equation}
Mỗi mẫu đầu vào có 21 checkpoint (input INT8, 4 pool output, GAP, 4 FC logit). Testbench
Questa nạp golden \texttt{.mem} và so sánh từng checkpoint. Kết quả lõi production:
\begin{itemize}
  \item 21/21 checkpoint PASS, $\max|\text{diff}| = 0$ LSB trên 15.312 phần tử so sánh
        (3 mẫu test).
  \item Latency xác định 5.216 chu kỳ cho mọi mẫu.
\end{itemize}

% =====================================================================
%  6. CROSS-DATASET
% =====================================================================
\section{Kiểm chứng cross-dataset với dữ liệu Georgia}

Để trả lời câu hỏi ``mô hình có tổng quát hóa ra ngoài dữ liệu huấn luyện không?'', dự án
đánh giá mô hình trực tiếp trên dữ liệu Georgia 12-lead (PhysioNet/CinC Challenge 2020, bộ
Emory G12EC). Đây là dữ liệu thu ở một quần thể và cơ sở y tế hoàn toàn khác, đóng vai
far-transfer cân bằng (balanced).

\subsection{Thiết lập đánh giá}
\begin{itemize}
  \item Bộ huấn luyện mở rộng: gộp Chapman và Ningbo (cùng chuẩn 12-lead/10s/500Hz, cùng
        taxonomy SNOMED 4 lớp) để mở rộng độ đa dạng quần thể và thiết bị thu, tăng độ bền của
        boundary giữa các lớp nhịp.
  \item Tiền xử lý: lấy lead II từ trường \texttt{val[1]} của file \texttt{.mat}, downsample
        500$\rightarrow$250 Hz, z-score chuẩn hóa, cửa sổ 2500 mẫu --- đồng nhất với pipeline
        Chapman.
  \item Ánh xạ nhãn: các SNOMED subclass của Georgia khớp đúng taxonomy 4 lớp của Chapman.
  \item Tập đánh giá Georgia: 5459 record, phân bố cân bằng AFIB 692 / GSVT 1192 / SB 1521 /
        SR 2054 --- khác hẳn PTB-XL (SR chiếm 84\%), nên là phép thử khắt khe về tổng quát hóa.
\end{itemize}

\subsection{Kết quả huấn luyện trên Chapman mở rộng đánh giá trên Georgia}

Mô hình huấn luyện trên Chapman + Ningbo (lượng tử hóa QAT-INT8, bit-exact với RTL) được
đánh giá zero-shot trực tiếp trên 5459 record Georgia:

\begin{table}[h]
\centering
\caption{Cross-dataset trên Georgia.}
\begin{tabular}{lccccc}
\toprule
\textbf{Chế độ} & \textbf{Acc} & \textbf{macro-P} & \textbf{macro-R} & \textbf{macro-F1} & \textbf{macro-AUC}\\
\midrule
Chapman/Ningbo $\rightarrow$ Georgia (INT8) & 0,9300 & 0,9142 & 0,9161 & 0,9151 & 0,9580\\
\bottomrule
\end{tabular}
\end{table}

Bảng per-class (train Chapman + Ningbo, INT8 $\rightarrow$ Georgia, 5459 record):

\begin{table}[h]
\centering
\caption{Per-class (train Chapman/Ningbo, INT8) trên Georgia.}
\begin{tabular}{lcccr}
\toprule
\textbf{Lớp} & \textbf{Precision} & \textbf{Recall} & \textbf{F1} & \textbf{Support}\\
\midrule
AFIB & 0,831 & 0,824 & 0,827 & 692\\
GSVT & 0,921 & 0,945 & 0,933 & 1192\\
SB & 0,954 & 0,962 & 0,958 & 1521\\
SR & 0,951 & 0,933 & 0,942 & 2054\\
macro & 0,914 & 0,916 & 0,915 & 5459\\
\bottomrule
\end{tabular}
\end{table}

Ma trận nhầm lẫn (hàng = nhãn thật, cột = dự đoán), INT8:

\begin{lstlisting}
             AFIB  GSVT    SB    SR
  AFIB  ->     570    35    11    76
  GSVT  ->      46  1127     7    12
  SB    ->      37    10  1464    10
  SR    ->      33    52    53  1916
\end{lstlisting}

\subsection{Phân tích}

Train Chapman/Ningbo, test Georgia (far-transfer balanced) --- cho một bức tranh tổng quát
hóa đầy đủ, trải từ phân phối gần đến phân phối xa; và cho thấy chiến lược mở rộng tập huấn
luyện sang dữ liệu cùng họ là hướng nâng độ bền hiệu quả mà không cần can thiệp vào dữ liệu
test.

% =====================================================================
%  7. DEMO
% =====================================================================
\section{Demo thiết kế trên FPGA}

Lõi được nạp lên kit DE10-Standard và điều khiển qua cầu JTAG-to-Avalon kết hợp công cụ
System Console của Quartus. Quy trình demo: nạp trọng số Chapman vào weight-RAM qua Avalon,
ghi mẫu ECG vào input SRAM, kích hoạt inference, đọc kết quả \texttt{result[1:0]}.
\begin{itemize}
  \item \textbf{Kết quả:} Phân loại tập test Chapman trên board đạt 94,27\% (1004/1065), tái
        lập kết quả mô phỏng $\sim$94\%. Chênh lệch là do test-subset/lần chạy, không phải sai
        số số học (datapath đã bit-exact).
  \item \textbf{Hạ tầng điều khiển:} JTAG-to-Avalon + System Console.
\end{itemize}
Demo này hoàn tất mục tiêu ``triển khai và đo đạc thực tế trên FPGA'': lõi không chỉ đúng
trong mô phỏng mà còn cho ra cùng kết quả khi chạy trên silicon thật.

% =====================================================================
%  8. TIẾN ĐỘ
% =====================================================================
\section{Báo cáo tiến độ hiện tại}

\begin{table}[h]
\centering
\caption{Trạng thái tiến độ theo Phase.}
\begin{tabularx}{\textwidth}{lXl}
\toprule
\textbf{Phase} & \textbf{Nội dung} & \textbf{Trạng thái}\\
\midrule
Software baseline & Prune (4,4,8,8), QAT-INT8 round-half-up, export hex, golden .mem & Hoàn thành\\
Cross-dataset & Georgia (far-balanced) & Hoàn thành\\
Hardware Design & ecg\_core.v + wrapper Avalon bus & Hoàn thành\\
Verification & Testbench 3 sample golden raw int8 & Hoàn thành\\
Synthesis & ALM 2201, DSP 28, Fmax 104,85 & Hoàn thành\\
Demo on-board & JTAG-to-Avalon DE10 94,27\% & JTAG done\\
\bottomrule
\end{tabularx}
\end{table}

\end{document}
